%! Carrying Out a Test for a Population Mean — Unit 7, Lesson 7.5 — Warm-Up
\documentclass[10pt]{article}
\usepackage{apstats-article}
\usepackage{apstats-boxes}
\newcommand{\TallMath}[1]{$\displaystyle #1\rule[-1.4em]{0pt}{3.2em}$}

\begin{document}

\pageheader{Unit 7, Lesson 7.5}{Warm-Up}
\namedateperiod

\vspace{0.15in}

\begin{notesbox}{Spiral Review --- Setting Up a $t$-Test (Lesson 7.4)}

A nutritionist claims the mean daily sodium intake of adults is 2,400 mg. A health researcher
believes the true mean is \emph{higher}. She takes a random sample of 35 adults.
$\bar x = 2{,}510$ mg and $s = 380$ mg. A histogram of the data shows a roughly symmetric
distribution with no outliers.

\textbf{1.} Name the inference method and give the degrees of freedom.

\quad Method: \blank{5cm} \quad $df = $ \blank{1.5cm}

\textbf{2.} State hypotheses.

\quad $H_0$: \blank{5cm} \quad $H_a$: \blank{5cm}

\textbf{3.} Verify both conditions.
\begin{itemize}
  \item Random: \writeline
  \item Normal: \writeline
\end{itemize}

\textbf{4.} Using the formula $t = \dfrac{\bar x - \mu_0}{s / \sqrt{n}}$, compute the
$t$-statistic. Show your work.

\vspace{0.15in}
$SE = \dfrac{s}{\sqrt{n}} = \dfrac{\blank{2cm}}{\sqrt{\blank{1.5cm}}} = $ \blank{2.5cm}
\hfill
$t = \dfrac{\blank{2cm} - \blank{2cm}}{\blank{2.5cm}} = $ \blank{2.5cm}

\end{notesbox}

\end{document}
